\subsection{Củng cố}

\begin{tuluyen}
Cho các chất: \ce{HCl}, \ce{NaOH}, \ce{NaCl}, \ce{Ca(OH)2}, \ce{H2SO4}. Phân loại thành acid, base và muối.
\dapso{Acid: \ce{HCl}, \ce{H2SO4}; base: \ce{NaOH}, \ce{Ca(OH)2}; muối: \ce{NaCl}.}
\end{tuluyen}

\begin{tuluyen}
Cho ba dung dịch có pH bằng 2, 7 và 11. Chọn dung dịch có tính acid mạnh hơn trong ba dung dịch và dung dịch có tính base.
\dapso{pH 2 có tính acid; pH 11 có tính base; pH 7 trung tính.}
\end{tuluyen}

\begin{tuluyen}
Viết phương trình hóa học và nêu hiện tượng khi cho magnesium vào dung dịch \ce{HCl} dư.
\dapso{\ce{Mg + 2HCl -> MgCl2 + H2 ^}; Mg tan dần, có khí không màu thoát ra.}
\end{tuluyen}

\begin{tuluyen}
Viết phương trình hóa học cho các phản ứng sau nếu có:
\begin{enumerate}
\item \ce{CuO} với \ce{HCl};
\item \ce{CO2} với \ce{Ca(OH)2};
\item \ce{NaCl} với \ce{KOH};
\item \ce{BaCl2} với \ce{Na2SO4}.
\end{enumerate}
\dapso{(a) \ce{CuO + 2HCl -> CuCl2 + H2O}; (b) \ce{CO2 + Ca(OH)2 -> CaCO3 v + H2O}; (c) không phản ứng; (d) \ce{BaCl2 + Na2SO4 -> BaSO4 v + 2NaCl}.}
\end{tuluyen}

\begin{tuluyen}
Nêu hiện tượng khi cho dung dịch \ce{NaOH} vào dung dịch \ce{CuSO4}. Viết phương trình hóa học.
\dapso{Xuất hiện kết tủa xanh \ce{Cu(OH)2}; \ce{CuSO4 + 2NaOH -> Cu(OH)2 v + Na2SO4}.}
\end{tuluyen}

\begin{tuluyen}
Viết các phương trình hóa học thực hiện chuỗi chuyển hóa:
\[
\ce{CaCO3 -> CaO -> Ca(OH)2 -> CaCO3}.
\]
\dapso{\ce{CaCO3 -> CaO + CO2}; \ce{CaO + H2O -> Ca(OH)2}; \ce{Ca(OH)2 + CO2 -> CaCO3 v + H2O}.}
\end{tuluyen}

\begin{tuluyen}
Có ba dung dịch không nhãn: \ce{HCl}, \ce{Na2CO3}, \ce{NaCl}. Chỉ dùng quỳ tím và dung dịch \ce{HCl} đã nhận biết được, hãy trình bày cách nhận biết cả ba dung dịch.
\dapso{Quỳ tím nhận \ce{HCl}; dùng \ce{HCl} nhận \ce{Na2CO3} qua khí \ce{CO2}; mẫu còn lại là \ce{NaCl}.}
\end{tuluyen}

\begin{tuluyen}
Trộn \SI{100}{\milli\liter} dung dịch \ce{HCl} \SI{0.20}{\mole\per\liter} với \SI{150}{\milli\liter} dung dịch \ce{NaOH} \SI{0.10}{\mole\per\liter}. Tính số mol chất dư.
\dapso{\ce{HCl} dư \SI{0.005}{\mole}.}
\end{tuluyen}

\begin{tuluyen}
Cho \SI{5.00}{\gram} \ce{CaCO3} phản ứng hoàn toàn với \ce{HCl} dư. Tính thể tích \ce{CO2} ở $V_m=\SI{24.8}{\liter\per\mole}$.
\dapso{\SI{1.24}{\liter}.}
\end{tuluyen}

\begin{tuluyen}
Trộn \SI{100}{\milli\liter} dung dịch \ce{CuSO4} \SI{0.30}{\mole\per\liter} với \SI{200}{\milli\liter} dung dịch \ce{NaOH} \SI{0.20}{\mole\per\liter}. Tính khối lượng kết tủa \ce{Cu(OH)2}. Cho $M=98$ g/mol.
\dapso{\SI{1.96}{\gram}.}
\end{tuluyen}

\begin{tuluyen}
Chuẩn độ \SI{25.0}{\milli\liter} dung dịch \ce{HCl} bằng dung dịch \ce{NaOH} \SI{0.100}{\mole\per\liter}. Ba lần chuẩn độ cần lần lượt 19,8 mL; 20,0 mL; 20,1 mL \ce{NaOH}. Lấy thể tích trung bình của ba lần, tính nồng độ mol của \ce{HCl}.
\dapso{$V_{tb}\approx\SI{19.97}{\milli\liter}$; $C_{\ce{HCl}}\approx\SI{0.0799}{\mole\per\liter}$.}
\end{tuluyen}

\begin{tuluyen}
Một mẫu đá vôi chứa \SI{80}{\percent} \ce{CaCO3} theo khối lượng. Cho \SI{250}{\gram} mẫu phản ứng hoàn toàn với \ce{HCl} dư. Tính khối lượng \ce{CaCl2} tạo thành. Cho $M_{\ce{CaCO3}}=100$, $M_{\ce{CaCl2}}=111$ g/mol.
\dapso{\SI{222}{\gram} \ce{CaCl2}.}
\end{tuluyen}

